\section{The Lorentz observation morphism and the full nine-dimensional fibre}
\label{hlob:sec:observation-morphism}

The rank-three Hermitian carrier and the four-real-dimensional Lorentz cone
are not identical objects.  The required next step is therefore to compute
their exact relation.  This section constructs a split positive observation
morphism, determines every one of its fibres, identifies the five retained
directions beyond the observed four-vector, and proves the complete cubic
coupling among the resulting Lorentz-vector, spinor, and scalar coordinates.
No direction is removed from the source calculation.

\subsection{The marked angular epimorphism and its exact splitting}
\label{hlob:subsec:marked-splitting}

Retain the angular lattice matrix and its primitive kernel vector:
\begin{equation}
 P=
 \begin{pmatrix}
  1&-1&0\\
  0& 1&-1
 \end{pmatrix},
 \qquad
 u=\begin{pmatrix}1\\1\\1\end{pmatrix},
 \qquad Pu=0.
 \label{hlob:eq:P-u}
\end{equation}
The standard Hermitian products in these displayed marked bases give
\begin{equation}
 C:=PP^*
 =\begin{pmatrix}2&-1\\-1&2\end{pmatrix},
 \qquad
 C^{-1}=\frac13\begin{pmatrix}2&1\\1&2\end{pmatrix},
 \qquad \det C=3.
 \label{hlob:eq:cartan}
\end{equation}
Define the corresponding orthogonal right inverse and the projection onto
the displayed kernel line by
\begin{equation}
 R:=P^*C^{-1}
 =\frac13
 \begin{pmatrix}
  2& 1\\
 -1& 1\\
 -1&-2
 \end{pmatrix},
 \qquad
 \Pi:=\frac{uu^*}{3},
 \qquad Q_u:=I_3-\Pi.
 \label{hlob:eq:R-Pi}
\end{equation}
The choice of \(R\) uses the displayed marking and its standard Hermitian
products.  All formulas below retain that choice explicitly.

\begin{lemma}[Exact splitting identities]
\label{hlob:lem:splitting-identities}
The matrices in \eqref{hlob:eq:P-u}--\eqref{hlob:eq:R-Pi} satisfy
\begin{equation}
 PR=I_2,\qquad u^*R=0,\qquad R^*R=C^{-1},\qquad
 RP=Q_u,\qquad \Pi R=0,\qquad P\Pi=0.
 \label{hlob:eq:splitting-identities}
\end{equation}
Moreover, with
\begin{equation}
 X=(R\;u)
 =\begin{pmatrix}
  2/3& 1/3&1\\
 -1/3& 1/3&1\\
 -1/3&-2/3&1
 \end{pmatrix},
 \label{hlob:eq:X}
\end{equation}
one has
\begin{equation}
 \det X=1,\qquad
 X^{-1}=\begin{pmatrix}P\\u^*/3\end{pmatrix}.
 \label{hlob:eq:X-inverse}
\end{equation}
\end{lemma}

\begin{proof}
Direct multiplication, with every denominator retained, gives
\[
 PR
 =\frac13
 \begin{pmatrix}1&-1&0\\0&1&-1\end{pmatrix}
 \begin{pmatrix}2&1\\-1&1\\-1&-2\end{pmatrix}
 =I_2
\]
and
\[
 RP
 =\frac13
 \begin{pmatrix}
  2&-1&-1\\
 -1& 2&-1\\
 -1&-1& 2
 \end{pmatrix}
 =I_3-\frac13
 \begin{pmatrix}1&1&1\\1&1&1\\1&1&1\end{pmatrix}
 =Q_u.
\]
The identities involving \(u\), \(R^*R\), and \(\Pi\) follow from the same
displayed matrices and from \(R=P^*(PP^*)^{-1}\).  Finally,
\[
 \det(3X)
 =\det\begin{pmatrix}2&1&3\\-1&1&3\\-1&-2&3\end{pmatrix}
 =27,
\]
so \(\det X=27/3^3=1\).  Both products
\[
 \begin{pmatrix}P\\u^*/3\end{pmatrix}(R\;u)
 \quad\text{and}\quad
 (R\;u)\begin{pmatrix}P\\u^*/3\end{pmatrix}
\]
are the identity by \eqref{hlob:eq:splitting-identities}; this proves
\eqref{hlob:eq:X-inverse}.
\end{proof}

\subsection{A positive observation morphism with all fibres retained}
\label{hlob:subsec:observation}

Put
\[
 \mathcal J_3=\operatorname{Herm}_3(\mathbb C),\qquad
 \mathcal J_2=\operatorname{Herm}_2(\mathbb C),\qquad
 \mathcal P_r=\{H\in\mathcal J_r:H\succeq0\}.
\]
Define real-linear maps
\begin{equation}
 \mathcal O_P:\mathcal J_3\longrightarrow\mathcal J_2,
 \qquad \mathcal O_P(H)=PHP^*,
 \label{hlob:eq:observation}
\end{equation}
and
\begin{equation}
 \jmath_P:\mathcal J_2\longrightarrow\mathcal J_3,
 \qquad \jmath_P(G)=RGR^*.
 \label{hlob:eq:section}
\end{equation}

\begin{theorem}[Split observation and the exact five-dimensional kernel]
\label{hlob:thm:split-observation}
The observation morphism is surjective and positive, and its displayed
section is a positive linear embedding:
\begin{equation}
 \mathcal O_P\jmath_P=\operatorname{id}_{\mathcal J_2},
 \qquad
 \jmath_P\mathcal O_P(H)=Q_uHQ_u.
 \label{hlob:eq:split-observation}
\end{equation}
Consequently
\begin{equation}
 \mathcal J_3
 =\operatorname{im}\jmath_P\oplus\ker\mathcal O_P,
 \qquad
 \dim_{\mathbb R}\operatorname{im}\jmath_P=4,
 \qquad
 \dim_{\mathbb R}\ker\mathcal O_P=5.
 \label{hlob:eq:dimension-split}
\end{equation}
The kernel is explicitly
\begin{equation}
 \ker\mathcal O_P
 =\{a\,uu^*+uv^*+vu^*:a\in\mathbb R,\ v\in u^\perp\}.
 \label{hlob:eq:linear-kernel}
\end{equation}
It is exactly the tangent space at \(H_0=uu^*\) to the nonzero rank-one
positive-semidefinite stratum
\begin{equation}
 \mathcal R_1=\{ww^*:w\in\mathbb C^3\setminus\{0\}\}
 \subset\mathcal P_3.
 \label{hlob:eq:rank-one-stratum}
\end{equation}
On the positive cone, the complete zero fibre is the retained ray
\begin{equation}
 \mathcal O_P^{-1}(0)\cap\mathcal P_3
 =\{\lambda uu^*: \lambda\geq0\}.
 \label{hlob:eq:positive-zero-fibre}
\end{equation}
\end{theorem}

\begin{proof}
Equations \eqref{hlob:eq:splitting-identities} give
\[
 \mathcal O_P(\jmath_P(G))=PRGR^*P^*=G
\]
and
\[
 \jmath_P(\mathcal O_P(H))=RPH P^*R^*=Q_uHQ_u.
\]
Both maps preserve positive semidefiniteness because they are congruences.
The first identity proves surjectivity, injectivity of \(\jmath_P\), and the
direct sum in \eqref{hlob:eq:dimension-split}.  Since
\(\dim_{\mathbb R}\mathcal J_3=9\) and
\(\dim_{\mathbb R}\mathcal J_2=4\), the dimensions follow.

The condition \(PHP^*=0\) is equivalent to \(Q_uHQ_u=0\).  Relative to
\(\mathbb C^3=u^\perp\oplus\mathbb Cu\), such a Hermitian operator has a
zero \(u^\perp\)-to-\(u^\perp\) block and therefore has the unique form
\eqref{hlob:eq:linear-kernel}.  The differential at \(u\) of
\(F(w)=ww^*\) is
\[
 dF_u(\dot w)=u\dot w^*+\dot w u^*.
\]
Writing \(\dot w=(a/2)u+v+i b u\), with
\(a,b\in\mathbb R\) and \(v\in u^\perp\), shows that its image is exactly
\eqref{hlob:eq:linear-kernel}; the imaginary radial direction \(ibu\) is
the one-real-dimensional kernel of \(dF_u\).  Thus the image has dimension
five and is the stated tangent space.

Finally let \(H\succeq0\) and \(PHP^*=0\).  For every \(\xi\in\mathbb C^2\),
\[
 0=\xi^*PHP^*\xi
   =\lVert H^{1/2}P^*\xi\rVert^2.
\]
Hence \(H^{1/2}\) annihilates
\(\operatorname{im}P^*=u^\perp\).  The range of \(H\) is contained in
\(\mathbb Cu\), and Hermitian positivity then gives
\(H=\lambda uu^*\) for one \(\lambda\geq0\).  The converse follows from
\(Pu=0\).
\end{proof}

The observation morphism is not a dimension count.  It factors through the
already proved two-column Gram carrier.  For \(H\in\mathcal P_3\), put
\begin{equation}
 \mathcal Z_P(H)=H^{1/2}P^*
 \in\operatorname{Mat}_{3\times2}(\mathbb C).
 \label{hlob:eq:gram-factor}
\end{equation}
Then, with no omitted scalar,
\begin{equation}
 \mathcal Z_P(H)^*\mathcal Z_P(H)=PHP^*=\mathcal O_P(H).
 \label{hlob:eq:gram-factor-identity}
\end{equation}

\subsection{Complete block coordinates and the cubic coupling}
\label{hlob:subsec:block-coordinates}

\begin{theorem}[Lorentz-vector--spinor--scalar coordinates]
\label{hlob:thm:block-coordinates}
Every \(H\in\mathcal J_3\) has a unique expression
\begin{equation}
 H=X
 \begin{pmatrix}
  G&\psi\\
  \psi^*&\lambda
 \end{pmatrix}X^*
 =RGR^*+R\psi u^*+u\psi^*R^*+\lambda uu^*,
 \label{hlob:eq:block-expression}
\end{equation}
where
\begin{equation}
 G=PHP^*,\qquad
 \psi=\frac13PHu,\qquad
 \lambda=\frac19u^*Hu,
 \label{hlob:eq:block-inverse}
\end{equation}
with \(G\in\mathcal J_2\), \(\psi\in\mathbb C^2\), and
\(\lambda\in\mathbb R\).  Thus the displayed marking gives the exact
real-vector-space decomposition
\begin{equation}
 \mathcal J_3\xrightarrow{\ \sim\ }
 \mathcal J_2\oplus\mathbb C^2\oplus\mathbb R,
 \qquad 9=4+4+1.
 \label{hlob:eq:four-four-one}
\end{equation}
For
\(G=\begin{psmallmatrix}a&b\\\bar b&d\end{psmallmatrix}\), retain
\begin{equation}
 \operatorname{adj}(G)
 =\begin{pmatrix}d&-b\\-\bar b&a\end{pmatrix}.
 \label{hlob:eq:adjugate}
\end{equation}
Then the complete determinant is
\begin{equation}
 \boxed{\det H
 =\lambda\det G-\psi^*\operatorname{adj}(G)\psi.}
 \label{hlob:eq:cubic-coupling}
\end{equation}

For \(G\succeq0\), let \(G^+\) be its Moore--Penrose inverse on
\(\operatorname{im}G\), extended by zero on \(\ker G\).  The complete
positive fibre is
\begin{equation}
 \begin{split}
 \mathcal O_P^{-1}(G)\cap\mathcal P_3
 =\biggl\{X
 \begin{pmatrix}G&\psi\\\psi^*&\lambda\end{pmatrix}X^*:
 &\ \psi\in\operatorname{im}G,\\[-1mm]
 &\ \lambda\geq\psi^*G^+\psi\biggr\}.
 \end{split}
 \label{hlob:eq:complete-positive-fibre}
\end{equation}
In particular,
\begin{equation}
 H\succ0
 \quad\Longleftrightarrow\quad
 G\succ0\ \text{ and }\ 
 \lambda>\psi^*G^{-1}\psi.
 \label{hlob:eq:positive-definite-condition}
\end{equation}
If \(G\) is singular and \(H\succeq0\), then \(H\) is singular.  If
\(G\succ0\), the boundary of its fibre is the graph
\(\lambda=\psi^*G^{-1}\psi\).
\end{theorem}

\begin{proof}
Equation \eqref{hlob:eq:X-inverse} gives
\[
 X^{-1}HX^{-*}
 =\begin{pmatrix}
   PHP^*&PHu/3\\
   u^*HP^*/3&u^*Hu/9
  \end{pmatrix},
\]
which proves uniqueness, \eqref{hlob:eq:block-expression}, and
\eqref{hlob:eq:block-inverse}.  Because \(\det X=1\), direct expansion of
the \(3\times3\) block determinant gives
\[
 \det H
 =\lambda(ad-|b|^2)
  -\begin{pmatrix}\bar\psi_1&\bar\psi_2\end{pmatrix}
   \begin{pmatrix}d&-b\\-\bar b&a\end{pmatrix}
   \begin{pmatrix}\psi_1\\\psi_2\end{pmatrix},
\]
which is \eqref{hlob:eq:cubic-coupling}.

Invertible congruence by \(X\) preserves positivity.  Suppose the block
matrix is positive semidefinite.  Its upper-left principal block gives
\(G\succeq0\).  If \(v\in\ker G\), positivity of the restriction to
\(\mathbb C(v,0)\oplus\mathbb C(0,1)\) forces \(v^*\psi=0\); hence
\(\psi\perp\ker G\), or equivalently
\(\psi\in\operatorname{im}G\).  Write \(\psi=G^{1/2}w\) with
\(w\in\operatorname{im}G\).  Completing the square gives
\[
 \begin{pmatrix}x\\s\end{pmatrix}^*
 \begin{pmatrix}G&\psi\\\psi^*&\lambda\end{pmatrix}
 \begin{pmatrix}x\\s\end{pmatrix}
 =\lVert G^{1/2}x+s w\rVert^2
  +|s|^2(\lambda-\lVert w\rVert^2).
\]
Here \(\lVert w\rVert^2=\psi^*G^+\psi\).  This proves necessity and
sufficiency of \eqref{hlob:eq:complete-positive-fibre}.  The strict case is
the ordinary Schur-complement criterion
\eqref{hlob:eq:positive-definite-condition}.

If \(G\) is singular, then \(\psi\in\operatorname{im}G\) and
\(\operatorname{adj}(G)\) has image in \(\ker G\).  Thus
\(\psi^*\operatorname{adj}(G)\psi=0=\det G\), so
\eqref{hlob:eq:cubic-coupling} gives \(\det H=0\).  If \(G\succ0\), the
same formula becomes
\[
 \det H=\det G\,(\lambda-\psi^*G^{-1}\psi),
\]
which proves the last assertion.
\end{proof}

\subsection{The embedded Lorentz group and the meaning of
\texorpdfstring{\(9=4+4+1\)}{9=4+4+1}}
\label{hlob:subsec:lorentz-action}

For \(g\in GL_2(\mathbb C)\), define
\begin{equation}
 T_P(g)=RgP+\Pi
 =X\begin{pmatrix}g&0\\0&1\end{pmatrix}X^{-1}.
 \label{hlob:eq:Tg}
\end{equation}

\begin{theorem}[Exact Lorentz-vector, Weyl-spinor, and scalar action]
\label{hlob:thm:lorentz-action}
The map \(T_P:GL_2(\mathbb C)\to GL_3(\mathbb C)\) is an injective group
homomorphism satisfying
\begin{equation}
 T_P(g)u=u,\qquad
 PT_P(g)=gP,\qquad
 T_P(g)R=Rg,\qquad
 \det T_P(g)=\det g.
 \label{hlob:eq:Tg-identities}
\end{equation}
Under congruence by \(T_P(g)\), the block coordinates transform exactly as
\begin{equation}
 (G,\psi,\lambda)\longmapsto
 (gGg^*,g\psi,\lambda).
 \label{hlob:eq:block-action}
\end{equation}
Consequently
\begin{equation}
 \mathcal O_P(T_P(g)HT_P(g)^*)
 =g\mathcal O_P(H)g^*,
 \qquad
 T_P(g)\jmath_P(G)T_P(g)^*=\jmath_P(gGg^*).
 \label{hlob:eq:equivariance}
\end{equation}

For \(g\in SL_2(\mathbb C)\), congruence on \(\mathcal J_2\) is the standard
double-cover action on Minkowski space.  In the retained coordinates
\begin{equation}
 G(t,x_1,x_2,x_3)=
 \begin{pmatrix}
  t+x_3&x_1-i x_2\\
  x_1+i x_2&t-x_3
 \end{pmatrix},
 \label{hlob:eq:pauli-coordinates}
\end{equation}
one has
\begin{equation}
 \det G=t^2-x_1^2-x_2^2-x_3^2,
 \qquad
 G\succeq0\Longleftrightarrow
 t\geq0,\quad t^2\geq x_1^2+x_2^2+x_3^2.
 \label{hlob:eq:lorentz-cone}
\end{equation}
The induced Lorentz-vector homomorphism
\begin{equation}
 \rho_{\mathrm{vec}}:SL_2(\mathbb C)\longrightarrow
 GL(\mathcal J_2),\qquad
 \rho_{\mathrm{vec}}(g)(G)=gGg^*,
 \label{hlob:eq:rho-vector}
\end{equation}
has kernel \(\{I_2,-I_2\}\) and image \(SO^+(1,3)\).  The embedding
\(T_P\), and hence the action on the full \(4+4+1\) carrier, is faithful.
Indeed, on the full nine-dimensional carrier
\(-I_2\) fixes \(G\) and \(\lambda\) but sends
\(\psi\mapsto-\psi\).  Thus the full carrier retains the spin double cover,
and \eqref{hlob:eq:four-four-one} is exactly the real representation sum
\begin{equation}
 \boxed{
 \text{Lorentz four-vector}\ \oplus\
 \text{defining complex spin representation }\mathbb C^2\ \oplus\
 \text{real scalar}.}
 \label{hlob:eq:representation-sum}
\end{equation}
With the convention that \(g\) acts on undotted spinor indices, the middle
summand is called a left Weyl spinor; reversing that convention reverses the
chirality name but not the displayed action \(\psi\mapsto g\psi\).
\end{theorem}

\begin{proof}
The second expression in \eqref{hlob:eq:Tg} and
\eqref{hlob:eq:X-inverse} prove the homomorphism, determinant, and
injectivity assertions.  Equations \eqref{hlob:eq:splitting-identities}
give the other identities directly.  Moreover,
\[
 T_P(g)X=X(g\oplus1),
\]
so congruence of \eqref{hlob:eq:block-expression} gives
\eqref{hlob:eq:block-action}.  The two equivariance identities follow.

For \(g\in SL_2(\mathbb C)\),
\[
 \det(gGg^*)=|\det g|^2\det G=\det G,
\]
and positive semidefiniteness is preserved.  Hence congruence defines a
homomorphism to the identity-component (proper orthochronous) Lorentz group of
\eqref{hlob:eq:lorentz-cone}.  If \(gGg^*=G\) for every Hermitian \(G\),
then \(gg^*=I_2\) by taking \(G=I_2\); the remaining equations say that \(g\)
commutes with every Hermitian matrix.  Thus \(g\) is scalar, and
\(\det g=1\) gives \(g=\pm I_2\).

The differential has zero kernel: if
\(AG+GA^*=0\) for all Hermitian \(G\), then \(A+A^*=0\) at \(G=I_2\), and
the remaining equation is \([A,G]=0\) for every Hermitian \(G\).  Hence \(A\)
is scalar; \(\operatorname{tr}A=0\) gives \(A=0\).  Both
\(SL_2(\mathbb C)\) and \(SO^+(1,3)\) have real dimension six, so the image
is an open subgroup of the connected group \(SO^+(1,3)\), hence the whole
group.  Finally, \eqref{hlob:eq:block-action} gives the action of the
central element and identifies the three standard real representations.
\end{proof}

The cubic in \eqref{hlob:eq:cubic-coupling} is invariant under the full
embedded \(SL_2(\mathbb C)\) action.  Indeed,
\begin{equation}
 \operatorname{adj}(gGg^*)=g^{-*}\operatorname{adj}(G)g^{-1}
 \qquad(g\in SL_2(\mathbb C)),
 \label{hlob:eq:adjugate-covariance}
\end{equation}
so
\((g\psi)^*\operatorname{adj}(gGg^*)(g\psi)
=\psi^*\operatorname{adj}(G)\psi\).
For the larger \(GL_2(\mathbb C)\) action it is instead a relative invariant:
\begin{equation}
 \det\!\left(T_P(g)HT_P(g)^*\right)
 =|\det g|^2\det H.
 \label{hlob:eq:GL2-relative-invariant}
\end{equation}
Thus no \(GL_2\)-invariance is being inferred from the \(SL_2\) formula.

\subsection{The Lorentz cone inside the rank-loss
\texorpdfstring{\(S^8\)}{S8} and its normal opening}
\label{hlob:subsec:S8}

Let
\begin{equation}
 e_2(K)=\frac12\bigl((\operatorname{tr}K)^2
                         -\operatorname{tr}(K^2)\bigr)
 \qquad(K\in\mathcal J_3).
 \label{hlob:eq:e2}
\end{equation}

\begin{theorem}[Proper Lorentz-cone section and determinant opening]
\label{hlob:thm:cone-section}
For every \(G\in\mathcal J_2\),
\begin{equation}
 \operatorname{rank}\jmath_P(G)=\operatorname{rank}G,
 \qquad
 \det\jmath_P(G)=0,
 \qquad
 e_2(\jmath_P(G))=\frac13\det G.
 \label{hlob:eq:e2-section}
\end{equation}
Hence \(\jmath_P\) maps the entire future Lorentz cone properly and
injectively into
\begin{equation}
 \partial\mathcal P_3
 =\{H\succeq0:\det H=0\}
 =\mathcal P_{3,\leq2}.
 \label{hlob:eq:rank-loss-boundary}
\end{equation}
It extends to an embedding of one-point compactifications
\begin{equation}
 \jmath_P^+:\mathcal P_2^+\cong B^4
 \lhook\joinrel\longrightarrow
 (\partial\mathcal P_3)^+\cong S^8.
 \label{hlob:eq:compactified-section}
\end{equation}
The timelike, nonzero null, and zero strata map respectively to rank two,
rank one, and rank zero.

For every real \(\varepsilon\), the affine slice through \(H_0=uu^*\)
satisfies the exact identity
\begin{equation}
 \boxed{
 \det\bigl(uu^*+\varepsilon\jmath_P(G)\bigr)
 =\varepsilon^2\det G.}
 \label{hlob:eq:determinant-opening}
\end{equation}
For \(\varepsilon>0\) and \(G\succeq0\), timelike \(G\) maps to the interior
of \(\mathcal P_3\), nonzero null \(G\) maps to its rank-two boundary, and
\(G=0\) maps to \(H_0\).  The restricted determinant Hessian is
\begin{equation}
 D^2(\det)_{H_0}
 \bigl(\jmath_P(G),\jmath_P(K)\bigr)
 =\operatorname{tr}G\operatorname{tr}K-\operatorname{tr}(GK),
 \label{hlob:eq:restricted-hessian}
\end{equation}
so
\begin{equation}
 \frac12D^2(\det)_{H_0}
 \bigl(\jmath_P(G),\jmath_P(G)\bigr)
 =\det G
 =t^2-x_1^2-x_2^2-x_3^2.
 \label{hlob:eq:hessian-lorentz}
\end{equation}

The surjection
\(\mathcal O_P:\partial\mathcal P_3\to\mathcal P_2\) is not proper and
therefore does not extend by sending infinity to infinity.  Thus
\eqref{hlob:eq:compactified-section} is an actual compactified embedding,
whereas the reverse compactified arrow is not supplied by
\(\mathcal O_P\).
\end{theorem}

\begin{proof}
Because \(R\) has complex rank two, congruence by \(R\) preserves the rank
of a \(2\times2\) Hermitian matrix.  Its image has rank at most two, so its
\(3\times3\) determinant vanishes.  When \(G\) is invertible, the product of
the two nonzero eigenvalues of \(RGR^*\) equals
\[
 \det(GR^*R)=\det G\det(C^{-1})=\frac13\det G.
\]
That product is \(e_2(RGR^*)\).  Both sides are polynomial in the four real
entries of \(G\), so the identity extends to singular \(G\), proving
\eqref{hlob:eq:e2-section}.  The map is a linear embedding between
finite-dimensional spaces, so its restriction to the closed cone is proper.
The compactification statement now follows from the previously proved
identifications
\(\mathcal P_2^+\cong B^4\) and
\((\partial\mathcal P_3)^+\cong S^8\).

Using \eqref{hlob:eq:X},
\[
 uu^*+\varepsilon\jmath_P(G)
 =X\begin{pmatrix}\varepsilon G&0\\0&1\end{pmatrix}X^*.
\]
Since \(\det X=1\), taking determinants proves
\eqref{hlob:eq:determinant-opening}.  The rank and positivity statements
follow from this same block matrix.  Polarizing
\(\det(sG+tK)\) gives
\[
 \det(sG+tK)
 =s^2\det G
  +st\bigl(\operatorname{tr}G\operatorname{tr}K
           -\operatorname{tr}(GK)\bigr)
  +t^2\det K,
\]
which proves \eqref{hlob:eq:restricted-hessian} and
\eqref{hlob:eq:hessian-lorentz}.

For nonproperness choose \(G_0=\operatorname{diag}(1,0)\) and
\[
 H_n=\jmath_P(G_0)+n uu^*,\qquad n\in\mathbb N.
\]
The two summands have orthogonal ranges, so \(H_n\succeq0\) has rank two
and lies in \(\partial\mathcal P_3\).  Its norm tends to infinity, while
\(\mathcal O_P(H_n)=G_0\) for every \(n\).  A continuous extension taking
the source point at infinity to the target point at infinity would therefore
send this convergent compactified sequence both to \(G_0\) and to infinity,
which is impossible.
\end{proof}

\subsection{The ten-dimensional determinant-phase carrier in the same coordinates}
\label{hlob:subsec:ten-dimensional}

The square-matrix \(SU(3)\)-orbit theorem already identifies its affine
quotient with
\begin{equation}
 \mathcal Y_3
 =\{(H,\omega):H\in\mathcal P_3,\ |\omega|^2=\det H\}.
 \label{hlob:eq:Y3}
\end{equation}
Substitution of \eqref{hlob:eq:block-expression} gives an exact
reformulation.

\begin{corollary}[Lorentz-spinor-scalar-phase form of the \(SU(3)\) quotient]
\label{hlob:cor:Y3-block-form}
There is a homeomorphism
\begin{equation}
 \mathcal Y_3\cong
 \left\{(G,\psi,\lambda,\omega):
 \begin{pmatrix}G&\psi\\\psi^*&\lambda\end{pmatrix}\succeq0,\quad
 |\omega|^2
 =\lambda\det G-\psi^*\operatorname{adj}(G)\psi
 \right\}.
 \label{hlob:eq:Y3-block-homeomorphism}
\end{equation}
The \(SL_2(\mathbb C)\)-action is
\begin{equation}
 (G,\psi,\lambda,\omega)\longmapsto
 (gGg^*,g\psi,\lambda,\omega).
 \label{hlob:eq:Y3-action}
\end{equation}
The complex determinant coordinate is invariant, and its zero locus is
exactly the rank-deficient cone:
\begin{equation}
 \omega=0\quad\Longleftrightarrow\quad\det H=0.
 \label{hlob:eq:omega-zero}
\end{equation}
After the proved one-point compactifications, this is the distinguished
\(S^8\) inside the ten-real-dimensional \(S^{10}\) carrier.  The ray
\(\{(G,\psi,\omega)=(0,0,0),\lambda\geq0\}\) is retained inside that
zero locus.
\end{corollary}

\begin{proof}
The block-coordinate map is an invertible real-linear congruence, so it
preserves positivity and is a homeomorphism.  Equation
\eqref{hlob:eq:cubic-coupling} gives the displayed scalar relation.
Equations \eqref{hlob:eq:block-action} and
\eqref{hlob:eq:adjugate-covariance} prove invariance.  The equivalence
\eqref{hlob:eq:omega-zero} follows directly from
\(|\omega|^2=\det H\).  The compactification and the identification of its
zero-determinant locus are the previously proved \(SU(3)\)-orbit and
rank-loss compactification theorems.
\end{proof}

This corollary makes the relation among the dimensions explicit.  The
nine-real-dimensional Hermitian carrier has the \(4+4+1\) decomposition in
\eqref{hlob:eq:four-four-one}.  Adding the complex determinant coordinate
and retaining its one-real-modulus equation gives \(9+2-1=10\) real
dimensions.  The determinant-zero locus removes the phase and imposes the
cubic boundary equation, leaving the eight-real-dimensional rank-loss cone
whose one-point compactification is \(S^8\).

\subsection{The holomorphic toric seed and its Lorentz-orbit completion}
\label{hlob:subsec:toric-seed}

The local four-branch chart has the already typed monomial map
\[
 q(z_0,z_1,z_2,z_3)=(z_0z_3,z_1,z_2).
\]
Define, without identifying angular and filling lattices,
\begin{equation}
 \kappa(w_0,w_1,w_2)=\operatorname{diag}(w_0,w_1w_2),
 \qquad
 K=\kappa q.
 \label{hlob:eq:kappa}
\end{equation}
Then
\begin{equation}
 K(z)=\operatorname{diag}(z_0z_3,z_1z_2),
 \qquad
 \boxed{\det K(z)=z_0z_1z_2z_3.}
 \label{hlob:eq:toric-seed}
\end{equation}
This is a holomorphic toric morphism into the diagonal complex matrices.
On the retained positive real locus put
\[
 a=z_0z_3,\qquad b=z_1z_2,\qquad
 p_0=\frac{a+b}{2},\qquad p_3=\frac{a-b}{2},\qquad p_1=p_2=0.
\]
Then \(K(z)=G(p)\succeq0\) and
\[
 p_0^2-p_3^2=ab=z_0z_1z_2z_3.
\]
Thus the toric products are exactly the two light-cone coordinates
\(p_0+p_3\) and \(p_0-p_3\).

For a retained \(\mu>0\), define the affine section
\begin{equation}
 \Psi_\mu(G)=\mu uu^*+\jmath_P(G).
 \label{hlob:eq:affine-section}
\end{equation}
Equation \eqref{hlob:eq:determinant-opening} gives
\begin{equation}
 \det\Psi_\mu(G)=\mu\det G,\qquad
 \Psi_\mu(G)\succeq0\Longleftrightarrow G\succeq0,\qquad
 \operatorname{rank}\Psi_\mu(G)=1+\operatorname{rank}G.
 \label{hlob:eq:affine-identities}
\end{equation}
Every \(G\succeq0\) has a spectral decomposition
\(G=v\operatorname{diag}(a,b)v^*\) with \(a,b\geq0\) and
\(v\in SU(2)\): a scalar phase changes the determinant of a chosen
unitary diagonalizer without changing its congruence action.  Hence
\begin{equation}
 \Theta_\mu:SL_2(\mathbb C)\times\mathbb R_{\geq0}^4
 \longrightarrow
 \Psi_\mu(\mathcal P_2),\qquad
 \Theta_\mu(g,z)=\Psi_\mu(gK(z)g^*)
 \label{hlob:eq:orbit-completion}
\end{equation}
is surjective.  It obeys
\begin{equation}
 \det\Theta_\mu(g,z)=\mu z_0z_1z_2z_3.
 \label{hlob:eq:orbit-completion-det}
\end{equation}
On \((\mathbb C^*)^4\), the complete auxiliary fibre of \(K\) is
\begin{equation}
 (\rho,\sigma)\cdot(z_0,z_1,z_2,z_3)
 =(\rho z_0,\sigma z_1,\sigma^{-1}z_2,\rho^{-1}z_3),
 \qquad(\rho,\sigma)\in(\mathbb C^*)^2.
 \label{hlob:eq:toric-fibre}
\end{equation}
Thus the local toric chart supplies a diagonal \(1+1\) seed, its
\(SL_2(\mathbb C)\)-orbit supplies all future causal directions, and the
two-complex-dimensional fibre in \eqref{hlob:eq:toric-fibre} is retained
as the exact redundancy of this parametrization.

\subsection{An exact colour-invariant Lorentz momentum and constrained dynamics}
\label{hlob:subsec:physical-translation}

The preceding results supply a pointwise carrier with commuting colour and
spin actions and a complete Lorentz-covariant orbit coordinate.  Let
\begin{equation}
 \mathscr V=\operatorname{Mat}_{3\times2}(\mathbb C),\qquad
 U\cdot Z=UZ\quad(U\in SU(3)),\qquad
 g\triangleright Z=Zg^*\quad(g\in SL_2(\mathbb C)).
 \label{hlob:eq:commuting-actions}
\end{equation}
These two left actions commute.  Retain positive conversion constants
\(\alpha>0\) and \(c>0\), with \(c\) carrying speed dimension, and define
\begin{equation}
 p_\alpha(Z)=\alpha Z^*Z\in\mathcal P_2,\qquad
 m_{\alpha,c}^2(Z)=\frac{\det p_\alpha(Z)}{c^2}.
 \label{hlob:eq:momentum-map}
\end{equation}
If \(p_\alpha\) is assigned momentum dimension, then
\([\alpha]=[p]/[Z]^2\) and
\([m_{\alpha,c}^2]=[p]^2/[c]^2=[\mathrm{mass}]^2\); no value of
\(\alpha\) or \(c\) is selected or suppressed.  Equivalently,
\(\det p_\alpha=m_{\alpha,c}^2c^2\), as required when the time component
of the Hermitian momentum matrix is \(E/c\).

\begin{theorem}[Complete colour-orbit momentum and its mass invariant]
\label{hlob:thm:colour-momentum}
The map \(p_\alpha\) is \(SU(3)\)-invariant, is covariant for the spin action,
and separates \(SU(3)\)-orbits:
\begin{equation}
 p_\alpha(UZ)=p_\alpha(Z),\qquad
 p_\alpha(g\triangleright Z)=g p_\alpha(Z)g^*.
 \label{hlob:eq:momentum-covariance}
\end{equation}
It induces a homeomorphism
\begin{equation}
 \mathscr V/SU(3)\xrightarrow{\ \sim\ }\mathcal P_2,
 \qquad [Z]\longmapsto p_\alpha(Z).
 \label{hlob:eq:momentum-quotient}
\end{equation}
Writing the rows of \(Z\) as three two-component covectors, one has
\begin{equation}
 \boxed{
 m_{\alpha,c}^2(Z)
 =\frac{\alpha^2}{c^2}\sum_{1\leq a<b\leq3}
 |Z_{a1}Z_{b2}-Z_{a2}Z_{b1}|^2.}
 \label{hlob:eq:mass-pluecker}
\end{equation}
It is invariant under both \(SU(3)\) and \(SL_2(\mathbb C)\), and
\begin{equation}
 m_{\alpha,c}^2(Z)>0\Longleftrightarrow\operatorname{rank}_{\mathbb C}Z=2,
 \qquad
 m_{\alpha,c}^2(Z)=0\Longleftrightarrow\operatorname{rank}_{\mathbb C}Z\leq1.
 \label{hlob:eq:mass-rank}
\end{equation}
For every rank-two \(Z_0\),
\begin{equation}
 m_{\alpha,c}^2(rZ_0)=r^4m_{\alpha,c}^2(Z_0)
 \qquad(r>0),
 \label{hlob:eq:mass-scaling}
\end{equation}
so the unrestricted kinematic carrier contains positive masses arbitrarily
close to zero.
\end{theorem}

\begin{proof}
The covariance equations follow directly from
\eqref{hlob:eq:commuting-actions}.  The orbit-separation and homeomorphism
statement is the exact two-column \(SU(3)\) Gram quotient, with multiplication
by the retained positive scalar \(\alpha\).  The \(2\times2\) Gram
determinant is the Cauchy--Binet identity
\[
 \det(Z^*Z)
 =\sum_{a<b}|Z_{a1}Z_{b2}-Z_{a2}Z_{b1}|^2,
\]
which proves \eqref{hlob:eq:mass-pluecker}.  Congruence by an
\(SL_2(\mathbb C)\) matrix preserves the determinant, while left \(SU(3)\)
preserves \(Z^*Z\).  The sum vanishes exactly when all \(2\times2\) minors
vanish, which is equivalent to rank at most one.  Finally
\(p_\alpha(rZ)=r^2p_\alpha(Z)\), and taking the \(2\times2\) determinant
gives \eqref{hlob:eq:mass-scaling}.
\end{proof}

For a retained mass \(m_0\geq0\), the constraint
\begin{equation}
 \mathscr C_{m_0,\alpha,c}
 =\{Z\in\mathscr V:\det p_\alpha(Z)=m_0^2c^2\}
 \label{hlob:eq:mass-shell-preimage}
\end{equation}
is invariant under both actions, and for \(m_0>0\) its \(SU(3)\)-orbit
space is exactly the future mass hyperboloid
\begin{equation}
 \{p\in\mathcal P_2:\det p=m_0^2c^2\}.
 \label{hlob:eq:mass-hyperboloid}
\end{equation}

Retain the Lorentz bilinear pairing
\begin{equation}
 \langle A,B\rangle_L
 =\frac12\bigl(\operatorname{tr}A\operatorname{tr}B
                    -\operatorname{tr}(AB)\bigr),
 \qquad
 \langle A,A\rangle_L=\det A.
 \label{hlob:eq:lorentz-pairing}
\end{equation}
For paths \(x:I\to\mathcal J_2\), \(Z:I\to\mathscr V\), and a real
multiplier \(e:I\to\mathbb R\), define the first-order constrained action
\begin{equation}
 \mathcal S_{m_0,\alpha,c}[x,Z,e]
 =\int_I
 \left(
  \langle p_\alpha(Z),\dot x\rangle_L
  -\frac{e}{2}\bigl(\det p_\alpha(Z)-m_0^2c^2\bigr)
 \right)d\tau.
 \label{hlob:eq:worldline-action}
\end{equation}
It is invariant under arbitrary \(\tau\)-dependent left \(SU(3)\) rotations
of \(Z\) and under the constant \(SL_2(\mathbb C)\) action
\begin{equation}
 x\longmapsto gxg^*,\qquad Z\longmapsto Zg^*.
 \label{hlob:eq:worldline-Lorentz-action}
\end{equation}
If \(x\) has length dimension and \(\tau\) retains its
chosen parameter dimension, then
\([e]=[x]/([\tau][p])\), so both summands in the integrand have dimension
\([p][x]/[\tau]\) and the action has dimension \([p][x]\).
Variation of \(e\) gives
\eqref{hlob:eq:mass-shell-preimage}; compactly supported variation of \(x\)
gives
\begin{equation}
 \frac{d}{d\tau}p_\alpha(Z)=0.
 \label{hlob:eq:momentum-conservation}
\end{equation}
For a variation \(\delta Z:I\to\mathscr V\), retain the exact differential
\begin{equation}
 \delta p_\alpha
 =\alpha\bigl(\delta Z^*Z+Z^*\delta Z\bigr).
 \label{hlob:eq:momentum-differential}
\end{equation}
The \(Z\)-Euler--Lagrange equation is the weak identity
\begin{equation}
 \left\langle
 \alpha\bigl(\delta Z^*Z+Z^*\delta Z\bigr),
 \dot x-e p_\alpha(Z)
 \right\rangle_L=0
 \quad\text{for every }\delta Z\in\mathscr V.
 \label{hlob:eq:Z-weak-equation}
\end{equation}
On the rank-two stratum this is equivalent to
\begin{equation}
 \dot x=e p_\alpha(Z).
 \label{hlob:eq:worldline-velocity}
\end{equation}
Indeed, for arbitrary \(A=A^*\in\mathcal J_2\), the variation
\begin{equation}
 \delta Z_A
 =\frac{1}{2\alpha}Z(Z^*Z)^{-1}A
 \label{hlob:eq:surjective-gram-variation}
\end{equation}
satisfies \(\delta p_\alpha=A\).  Hence the differential of the Gram map is
surjective there, and nondegeneracy of \(\langle\ ,\ \rangle_L\) proves
\eqref{hlob:eq:worldline-velocity}.  On the rank-at-most-one stratum the
differential need not be surjective, so the complete equation is
\eqref{hlob:eq:Z-weak-equation}; it is not replaced by the rank-two formula.
The invariance of \eqref{hlob:eq:lorentz-pairing} follows by polarizing
\(\det(gAg^*)=\det A\).  The multiplier and position equations follow by
direct variation and integration by parts, while
\eqref{hlob:eq:Z-weak-equation} follows from
\(D(\det)_p[A]=2\langle p,A\rangle_L\).

Equation \eqref{hlob:eq:worldline-action} is a complete constrained
kinematic dynamics with an independently retained mass parameter.  It does
not dynamically select \(m_0>0\): choosing \(m_0=0\) is equally compatible,
and \eqref{hlob:eq:mass-scaling} proves that the unrestrained carrier has no
positive kinematic threshold.  The next physical construction is thereby
specified: introduce a field-theoretic kinetic operator and interaction on
the full \((G,\psi,\lambda,\omega)\) carrier, derive rather than insert the
selected spectral sector, and compare its gauge-invariant spectrum with the
measured regime.  The exact morphisms and representations proved above are
the data that such a construction must preserve.
